\section{Introduction}

When the Census Bureau counts the population of a county, it counts every person physically present at the time of enumeration---including the 2.1 million people housed in state prisons, federal penitentiaries, county jails, and other correctional facilities. The Bureau's ``usual residence'' rule defines a person's home as the place where they live and sleep most of the time, and for incarcerated people, that is the facility, not their pre-incarceration address.

This counting convention has a redistributive consequence that was invisible when mass incarceration was small and geographically diffuse. As the carceral population grew from approximately 500,000 in 1980 to over 2 million by 2008, and as prisons became concentrated in rural areas (partly due to economic development strategies in struggling rural communities), the political effect became measurable: rural districts containing large prisons receive population credit---and therefore political representation---for constituents who cannot vote, who typically cannot access local government services at the prison location, and who overwhelmingly originate from urban areas hundreds of miles away.

The Prison Policy Initiative has documented this effect comprehensively for state legislative redistricting~\citep{wagner2002importing, wagner2004prison}. Our contribution is to quantify the effect for congressional redistricting and to analyze how algorithmic redistricting systems handle---and could handle---prison-population corrections.

\subsection{Scope and Definitions}

We use the following terms throughout:
\begin{itemize}
    \item \textbf{Prison gerrymandering}: The distortion of district boundaries caused by counting incarcerated people at their facility location rather than their home address.
    \item \textbf{Rural inflation}: The percentage increase in a county's or district's Census population attributable to incarcerated non-residents.
    \item \textbf{Residence-address correction}: The reallocation of prisoner population counts from the facility location to the prisoner's last home address, using administrative records or statistical models.
\end{itemize}

\subsection{Research Questions}

\begin{enumerate}
    \item How large is the rural inflation effect in prison-dense counties, and which states are most affected?
    \item What is the demonstrated experience of New York (reformed) versus Pennsylvania and Texas (not reformed) in managing prison gerrymandering?
    \item How would BISECT's algorithmic redistricting change under a prison-adjusted balance metric, and what implementation challenges does such a correction present?
\end{enumerate}

\subsection{Contributions}

\begin{itemize}
    \item Quantification of prison-county population inflation (8--22\%) across the 50 largest prison-dense counties using 2020 Census and DOJ facility data
    \item Three-state comparative case study: New York (reformed 2010), Pennsylvania (unreformed), Texas (unreformed, largest state prison system)
    \item Description of BISECT's planned \texttt{--balance-metric prison\_adjusted} preprocessing mode and its methodological requirements
    \item Legal analysis connecting prison gerrymandering to the VRA's vote-dilution standard
    \item Racial equity analysis demonstrating that prison gerrymandering transfers representation from Black and Hispanic communities to white communities
\end{itemize}
